\subsection{What the original measured}

The dissertation's conclusion was that echo state networks cannot learn the
variational level set, and that the leaking rate and spectral radius are
responsible. The first half is not supported by the experiments reported here:
under every configuration in which the models train at all, the echo state
network is within seed noise of the trained recurrent networks. The second half
is a claim about a parameter sensitivity that this study cannot confirm either
(Table~\ref{tab:blockA} and the sweep): the response surface over spectral
radius and leak is not systematic at the resolution a single seed per cell can
resolve, and the differences it shows are smaller than the seed to seed spread
of a fixed configuration.

What the original did measure, we think, is three properties of its own setup:
an optimiser that prevents training, a readout that destroys the geometry the
task is about, and a horizon at which the identity function is close to
optimal. None of these is a property of reservoir computing.

\subsection{Why the head matters so much here}

The target mask differs from the input mask in a small band around the moving
front. A decoder with a skip connection can express "copy, then edit the band";
a globally pooled vector cannot, because every pixel of the output must be
reconstructed from a representation that has already discarded position. This
is why the head changes results by more than the choice of recurrent cell does,
and why a study comparing cells has to control for it.

\subsection{What the reservoirs are actually worth}

On accuracy, nothing separates them from trained networks on this task. On
cost, the picture depends entirely on how the readout is fitted. Trained by
gradient descent, as the original did, a reservoir costs what a gated network
costs and is therefore pointless. Fitted in closed form, it is more than an
order of magnitude cheaper and gives up a substantial amount of accuracy.

The spiking liquid adds a further consideration the analogue reservoir does
not: its state is a filtered spike train with a measured firing rate of about a
quarter, which is the quantity that would matter on neuromorphic hardware where
cost is counted in spikes rather than in multiplications. On this task it
matches the analogue reservoir when its readout is trained and falls behind it
when the readout is solved.

\subsection{A methodological point that outlives this task}

A slowly evolving partial differential equation is an adversarial case for
whole image overlap metrics. Any metric dominated by the static part of the
signal will rank a model that predicts stasis above a model that predicts
motion imperfectly, and both above nothing. Three cheap habits avoid this:
report the identity baseline, restrict the metric to the region that changes,
and evaluate at a horizon long enough that the identity baseline fails. The
first alone would have caught the problem in the original study.
